\chapter{总结与展望}
\vspace{7pt}

\section{全文工作总结}

本文针对软件工程需求分析阶段中手动UML建模效率低下、语义偏差大以及大语言模型单次生成容易产生逻辑幻觉等工程痛点，设计并实现了一套基于多智能体协作与检索增强生成（RAG）技术的自动化用例建模系统。在研究过程中，本文摒弃了简单依赖大语言模型单次提示词输出的传统思路，转而构建了以LangGraph为核心的智能体状态机编排架构。通过将复杂的建模任务解构为实体提取、关系推理、属性分析等一系列专业化子任务，并利用共享状态池实现各Agent间的上下文传递，有效降低了模型在处理复杂逻辑时的认知负载。与此同时，系统引入了人机的交互机制，在关键建模节点设置断点，赋予用户对中间结果的审核与修正权，从而在保障建模效率的同时显著提升了结果的可靠性。

在解决长文档处理难题方面，本文探索了RAG技术在需求工程领域的应用落地。通过构建基于ChromaDB的语义向量库，实现了对大规模需求文档的智能分块与关联上下文检索，缓解了大语言模型在长程语义依赖下的信息遗忘与幻觉问题。针对后端生成的PlantUML代码，本文提出了一套包含多余空白行清洗、特殊字符转义以及逻辑引用校验的优化算法，消补了大模型模糊输出与建模引擎严苛语法要求之间的鸿沟，保证了生成代码的一次渲染成功率。

系统评测结果表明，本文提出的多Agent分步提取方案在类图、用例图和时序图的生成质量上均系统性地优于传统的单Agent生成方案。综合F1值从0.50至0.55的初级水平提升到了0.63至0.70的中等实用区间，特别是在逻辑关联推理的严谨性方面表现出显著优势。综上所述，本系统成功将先进的大语言模型推理能力与严谨的软件工程建模规范相结合，为实现需求分析到设计模型的自动化转化提供了一种解决方案。

\section{研究局限性分析}

尽管本系统在多Agent协同与工程化渲染方面取得了一定进展，但在实验测试与实际运行中仍显现出一些局限性。最为核心的挑战源于自动化建模效果对原始需求文档撰写质量的强依赖性。正如实验中所观察到的，当原始文本存在术语不一致、逻辑交织或显式动词缺失时，系统容易陷入“垃圾进、垃圾出”的困境，导致实体识别和关系判定的准确率大幅下降。这说明目前的系统尚缺乏对低质量、非结构化需求文本的深度重构与纠错能力。

在细粒度语义提取层面，系统表现尚不够稳健。类图的属性与方法提取、时序图的消息参数编排等任务对上下文细节的要求极高，而当前基于通用大模型的Agent在处理这些微观建模要素时，召回率往往处于较低水平，难以捕获需求文档中隐含的局部状态变量。此外，RAG模块在跨模块一致性方面的改善效果依然有限。目前的检索机制主要基于语义相似度，但在处理具有高度抽象关系的跨模块逻辑时，单纯的向量匹配有时无法准确召回关键的功能依赖信息，甚至可能引入无关噪声，干扰模型的判断。最后，系统在应对极端复杂、具有上百个实体的超大型项目时，状态池的维护开销与Agent间的通信复杂性会显著增加，目前的架构在扩展性与长序列推理稳定性上仍有待进一步验证。

\section{未来研究方向}

站在大语言模型赋能软件工程的技术前沿，本研究未来的演进方向将聚焦于模型专用化、多模态融合以及深度协作三个维度。首先，随着开源模型社区的发展，利用领域特定语料对大语言模型进行微调（Fine-tuning）将成为提升建模准确率的关键。近期的相关研究已经证明了这一方向的可行性，例如有学者通过构建包含多领域文字需求与对应 PlantUML 代码的数据集，对 Qwen2 等开源大模型进行微调和量化部署，有效提升了模型在生成 UML 类图时的质量\cite{li2025domainuml}。基于这一思路，未来计划针对UML语义规范和软件设计模式构建规模更大、质量更高的指令微调数据集，训练专门适配多类型UML图表协同建模任务的轻量化模型，以进一步弥补通用模型在细粒度属性提取与复杂逻辑推演上的短板。

多模态需求理解是另一个极具潜力的研究点。实际工程中，需求往往以“文字+原型图+流程图”的形式并存。未来可以将系统现有的文本分析引擎与视觉感知模型进行深层耦合，实现从手绘草图或业务流程图中辅助提取建模信息，从而提供更全面的信息源支持。此外，检索增强机制也需向知识图谱（GraphRAG）方向演进，通过构建需求实体间的显式关系图谱，替代单一的向量相似度检索，以从根本上解决跨模块建模的一致性难题。

在人机协作与迭代演进方面，未来的研究将致力于构建更加智能的增量式建模环境。大语言模型在处理持续迭代的业务需求时已展现出支持全生命周期建模的巨大潜力\cite{xie2026iterative}。因此，系统不仅应支持人工修正，还应学习用户的修改习惯，实现Agent提示词的动态自我优化。通过引入强化学习机制，让智能体在与人类分析师的反复交互中不断对齐专业建模偏好。最终目标是向“语义驱动、结构智能、人机协同优化”的新型协同建模范式迈进\cite{song2025umlai}，使系统演化为一个具备深厚工程知识沉淀、能够与人类专家深度共创的智能化软件建模助理。